% !PRINTED none
% part title: Abschnitt-opening page, no printed page number visible in the scan (chapter divider + opening section text; expected sequence puts this at p.548, but no digit is visible anywhere on the page image).
\begin{center}
{\large\gesperrt{Zweiundzwanzigster Abschnitt.}}

{\Large\textbf{Cayleys Entwickelung der Modulfunktionen.}}

\rule{0.3\textwidth}{0.4pt}

\textbf{\S\,145. Grenzwerte für $s = 1$.}
\end{center}

In diesem Abschnitt soll eine funktionentheoretische Anwendung der Grenzformel gegeben werden. Es bedeutet also hier $\omega$ nicht eine quadratische Irrationalzahl, sondern eine Variable mit positiv imaginärem Bestandteil. Die Modulfunktionen gehören zu den Funktionen mit \gesperrt{natürlichen Grenzen}, d.\,h. wenn man sich der Grenze der Konvergenz nähert, so liegen auf dieser Grenze, hier also auf der Achse der reellen $\omega$, die singulären Punkte überall dicht, so daß man diese Funktionen über diese Grenze hinaus nicht analytisch fortsetzen kann. \gesperrt{Cayley} hat in seinen letzten Untersuchungen Entwickelungen der Modulfunktionen gegeben, die darum merkwürdig sind, weil sie das Verhalten der Funktionen bei der Annäherung an die Grenze augenfällig machen. Der Schlüssel zu diesen Entwickelungen ist die Grenzformel § 141, (27)$^{1)}$.

Wenn man in dieser Formel:

\weq{1}{\begin{aligned}
&\mathop{\mathrm{Lim}}\limits_{s=1} \sum^{x,y} \frac{1}{(Ax^2+2Bxy+Cy^2)^s} - \frac{\pi}{(s-1)\sqrt m}\\
&= -\frac{2\pi\Gamma'(1)}{\sqrt m} + \frac{\pi}{\sqrt m}\log\frac{A}{4m} - \frac{2\pi}{\sqrt m}\log \eta(\omega_1),\ \eta(\omega_2),
\end{aligned}}
\[
\omega_1 = \frac{B+i\sqrt m}{A}, \quad \omega_2 = \frac{-B+i\sqrt m}{A},
\]

\weq{2}{\omega_1 = \alpha+\beta i, \quad \omega_2 = -\alpha+\beta i, \quad \beta > 0}

\begingroup\renewcommand{\thefootnote}{1)}\footnotetext[1]{Die erste \gesperrt{Cayley}sche Arbeit findet sich in dem Comptes Rendus der Pariser Akademie von 1893, Bd.\,161. Weiteres in einem Briefwechsel mit dem Verfasser dieses Buches, der im 47.\ Bande der mathematischen Annalen veröffentlicht ist.}\endgroup
%CONT>>
% [sic] Gedruckt steht Bd.161; die Jahresangabe1893 entspricht Bd.116 im Katalog der Bayerischen Staatsbibliothek. Quellenlesart unveraendert.
% [sic] Das Komma zwischen den eta-Faktoren in(1) steht so im Druck; (3) auf der Folgeseite bestaetigt das gemeinte Produkt.
